This thesis set out to explore whether or not the clause evaluation speedup achieved on the small NoisyXOR problem in a preceding project thesis \textcolor{red}{cite jonas den kule} carries over to MNIST scale, where the TM model far exceeds the on-chip memory of the NEORV32 processor and must reside in external memory (DDR2 SDRAM) at run time. To answer this, a multi-class MNIST TM with 2000 clauses per class was deployed on a Nexys A7 board featuring a Xilinx Artix-7 FPGA, and three inference implementations, a baseline software version and two CFU designs, were evaluated under a per-chunk and a per-class model load strategy.

The measurements lead to a clear conclusion. All six configurations achieved similar classification accuracy, confirming that every implementation evaluates the same model. The custom instructions, however, produced no meaningful speedup, with the simpler combinatorial CFU design improved inference speed by less than 1\%, and the more complex stateful CFU design performing worse than the baseline implementation. This stands in sharp contrast to the $4.42\times$ speedup reported at NoisyXOR scale.

The reason for this is a shift in the where the bottleneck of inference lies. At MNIST scale, moving the model out of DDR2, rather than evaluating clauses, takes up a majority of inference time.  In the per-class strategy the model preload alone accounts for approximately 91\% of the total, and the effective transfer rate of around 5.1\,MB/s reaches only a fraction of the maximum bandwidth the DDR2 device, indicating that the per-transaction overhead on single-word reads, not the memory device itself, is the limiting factor. Clause scoring contributes too small a share of inference time for any clause-level instruction to influence the result. The bottleneck is therefore not fixed by the algorithm but moves with the scale of the model, and at a realistic problem size it is the memory hierarchy rather than the computational power that limits performance.

These findings reframe processor-level TM acceleration as a memory problem rather than a compute problem once the model no longer fits on-chip. This opens up new directions for future work to take, like accelerating transfer rate through DDR2 burst transfers, a hybrid load strategy that preserves early exit, or a wider processor-external accelerator. 

It is also worth looking at other TM variants such as the Coalesced and Convolutional TM, which achieve better accuracy with far fewer clauses, in turn making the models far smaller, in addition to them having heavier computation costs per-clause. These factors might shift the bottleneck away from memory accessing, and custom instructions might still be effective for acceleration. On-chip training is feasible, but likely only as a reduced-scale proof of concept. The custom instructions themselves were shown to be cheap and correct; the limiting factor lies not in their design but in where the performance cost of inference at scale actually resides.